\par\Needspace{6\baselineskip}\medskip
\hypertarget{native-209-article155}{}\pdfbookmark[2]{155.}{article-155}
\begin{center}155.\end{center}
To show more clearly the precision of the elements we have found, we shall compute the middle position from them. For Oct.\ $17{,}415\,011$ the mean anomaly is found to be $=332^{0}\,28'\,54''{,}76$, whence the true anomaly $315^{0}\,1'\,23''{,}02$ and $\log r'=0{,}325\,9877$ (see \SourceLeaf{210}the examples in Articles 13, 14). This true anomaly should equal the true anomaly at the first observation increased by the angle $2f''$, or that at the third observation diminished by the angle $2f$, i.\ e.\ $=315^{0}\,1'\,22''{,}98$. The logarithm of the radius vector, in turn, should be $=0{,}325\,9878$: the differences may be disregarded. If the computation for the middle observation is continued as far as the geocentric position, the resulting numbers differ from the observation by only a few hundredths of a second of arc (Article 63). Such differences are, as it were, absorbed by the unavoidable errors arising from the limited precision of the tables.

We have treated the preceding example with the utmost precision to show how easily a solution of the highest accuracy can be obtained by our method. In practice it will rarely be necessary to follow this example so meticulously: it will generally suffice to use \textit{six} decimal digits throughout. In our example the second hypothesis would already have afforded no less precision, and the first would have afforded more than enough. We think readers will welcome a comparison between the elements obtained from the third hypothesis and those that would result if the second or even the first hypothesis were used in the same way for the same purpose. We present these three sets of elements in the following table:
\begin{center}
\begingroup\footnotesize\setlength{\tabcolsep}{2pt}\renewcommand{\arraystretch}{1.2}
\begin{tabular}{|p{1.35in}|r|r|r|}\hline
& from hypothesis III & from hypothesis II & from hypothesis I\\\hline
{\raggedright Epoch of mean long.\ 1805\par} & $41^{0}\,52'\,21''{,}68$ & $41^{0}\,52'\,18''{,}40$ & $42^{0}\,12'\,37''{,}83$\\
mean daily motion & $824''{,}7990$ & $824''{,}7978$ & $823''{,}5026$\\
perihelion & $52^{0}\,18'\,9''{,}30$ & $52^{0}\,18'\,6''{,}66$ & $52^{0}\,41'\,9''{,}81$\\
$\varphi$ & $14\quad12\quad1{,}87$ & $14\quad11\quad59{,}94$ & $14\quad24\quad27{,}49$\\
log.\ of semimajor axis & $0{,}422\,4389$ & $0{,}422\,4394$ & $0{,}422\,8944$\\
ascending node & $171^{0}\,7'\,48''{,}73$ & $171^{0}\,7'\,49''{,}15$ & $171^{0}\,5'\,48''{,}86$\\
inclination of orbit & $13\quad6\quad44{,}10$ & $13\quad6\quad45{,}12$ & $13\quad2\quad37{,}50$\\\hline
\end{tabular}\endgroup
\end{center}

Computing the heliocentric position in the orbit for the middle observation with the second set of elements gives an error in the logarithm of the radius vector $=0$ and an error in the longitude in the orbit $=0''{,}03$. Computing that position with the set derived from the first hypothesis gives an error in the logarithm of the radius vector $=0{,}000\,0002$ and an error in the longitude in the orbit $=1''{,}31$. Continuing the computation to the geocentric position gives
\SourceLeaf{211}
\begin{center}
\begingroup\small\setlength{\tabcolsep}{4pt}\renewcommand{\arraystretch}{1.2}
\begin{tabular}{|l|r|r|}\hline
& from hypothesis II & from hypothesis I\\\hline
Geocentric longitude & $352^{0}\,34'\,22''{,}26$ & $352^{0}\,34'\,19''{,}97$\\
error & $0{,}14$ & $2{,}15$\\
geocentric latitude & $-6\quad21\quad55{,}06$ & $-6\quad21\quad54{,}47$\\
error & $0{,}01$ & $0{,}60$\\\hline
\end{tabular}\endgroup
\end{center}
